\scanpage{3}
UNIVERSITE

DES SCIENCES ET TECHNIQUES

DU LANGUEDOC

\texttt{------------}

Service de la Scolarité

MA/GJ

Diffusion : Liste A

N° : 5982

Montpellier, le 12 Novembre 1980

AVIS DE SOUTENANCE D'UNE THESE DE DOCTEUR INGENIEUR

CHIMIE ORGANIQUE PHYSIQUE

\texttt{-------------------------}

\uline{Nom et prénom du candidat} : MAUBERT Claude

\uline{Titre de la thèse} : ETUDES CINETIQUES DES TELOMERISATIONS AMORCEES PAR LES SELS DE METAUX DE TRANSITION ET LES GENERATEURS DE RADICAUX LIBRES.

\uline{Date, lieu et heure de soutenance prévus} : 18 NOVEMBRE 1980 à 10 Heures 30 mn

E.N.S.C.M.

\uline{Nom du Directeur de thèse} : Monsieur le Professeur Y. PIETRASANTA

\uline{Désignation du laboratoire où la thèse est préparée} : Laboratoire de Chimie Appliquée

E.N.S.C.M. 8 Rue de l'Ecole Normale, 34075 MONTPELLIER CEDEX.

\uline{Résumé de la thèse} :

Cette thèse est axée sur la détermination des paramètres permettant de contrôler les réactions de télomérisation. Le cas des télomérisations par catalyse rédox des acrylates de méthyle et d'éthyle avec le tétrachlorure de carbone a d'abord été étudié. Il est apparu que les chlorures ferriques et cuivriques étaient les réactifs prépondérants pour le contrôle des degrés de polymérisation. On a ensuite mis en évidence une diminution de l'importance du transfert lorsque la longueur de la chaîne augmentait. Le dernier volet de cette étude concerne la télomé-\\
risation de l'acétate de vinyle avec le tétrachloro 2,4,4,4, butanoate de méthyle amorcée par l'azobisisobutyronitile. La constante de transfert au télogène a été calculée ; on a également évalué le pourcentage des chaînes macromoléculaires comportant une molécule de télogène en extrémité.

\uline{Durée approximative des recherches ayant abouti à la thèse} : 3 ans

\uline{Composition du jury proposé} :
\begin{tabular}{ll}
MM.&REMPP P.\\
& BOUTEVIN B.\\
& DEISS W.J.\\
& GENESTE P.\\
& PIETRASANTA Y.
\end{tabular}

P. le Secrétaire Général,

L'Attachée d'Administration Universitaire,

\uline{\unsure{P. Gentil|A. Gentil}}

\uline{A. GENTIL}
\mspage{2}
